%% ma: L12-B14-0112
%% lop: 12
%% chuong: 5
%% bai: 14
%% dang: 12.14.32
%% loai: TN4
%% muc: VD
%% dap_an: "D"
%% nguon: "Ảnh giáo viên gửi 07/11/2026 (Tư Duy Mở - tuduymo.com, đề tự luyện Phương trình mặt phẳng, đề khoảng cách), Câu 8"
%% kien_thuc: "Bài 14 (hai mặt phẳng song song, khoảng cách từ một điểm đến mặt phẳng)"
%% trang_thai: chinh-thuc
%% ghi_chu: "Tài liệu gốc không kèm đáp án; đáp án do Claude giải và kiểm bằng sympy. Đề gốc cho khoảng cách bằng $2$: khi đó $D=-1$ (trùng $(P)$, loại) hoặc $D=-13$ ($b=13$), không có phương án đúng (ý đề gốc chọn A ứng với mặt phẳng trùng $(P)$). Đã sửa khoảng cách thành $1$, giữ nguyên phương án; đáp án D. Dạng đã duyệt 07/11/2026."
\begin{ex}
Trong không gian $Oxyz$, cho mặt phẳng $(P)\colon x-2y+2z-1=0$ và điểm $A(1;0;3)$. Mặt phẳng $(Q)$ song song với $(P)$ và cách điểm $A$ một khoảng bằng $1$. Mặt phẳng $(Q)$ cắt trục $Ox$ tại điểm $B(b;0;0)$ có $b<5$. Giá trị của $b$ bằng
\choice{$1$}{$2$}{$3$}{\True $4$}
\loigiai{$(Q)\colon x-2y+2z+D=0$ ($D\ne-1$). $\dfrac{|1+6+D|}{3}=1\Leftrightarrow D=-4$ hoặc $D=-10$. $(Q)$ cắt $Ox$ tại $b=-D$, nên $b=4$ hoặc $b=10$; $b<5\Rightarrow b=4$. Chọn D.}
\end{ex}
